% ch7_d 第7章 §7.4.2尾–§7.4.5 (书页 314–326 / p329–p341)
%JOIN%
由电子算符 $\Psi^\dagger$ 产生。上述边缘激发理论是用一维密度算符 $\rho(x)$ 来表述的。因此，中心任务是把电子算符用密度算符表示出来。边缘上的电子算符产生一个局域电荷，故应满足
\begin{equation*}
[\rho(x),\Psi^\dagger(x')]=\delta(x-x')\Psi^\dagger(x')
\tag{7.4.6}
\end{equation*}
由于 $\rho$ 满足 K--M 代数 (7.4.5)，可以证明
\begin{equation*}
[\rho(x_1),\rho(x_2)]=\mathrm{i}\frac{\nu}{2\pi}\delta'(x_1-x_2)
\end{equation*}
或者
\begin{equation*}
[\rho(x_1),\phi(x_2)]=-\mathrm{i}\nu\delta(x_1-x_2)
\end{equation*}
其中 $\phi$ 由 $\rho=\frac{1}{2\pi}\partial_x\phi$ 给出。我们看到，$\rho(x)$ 可以视为 $\phi$ 的泛函导数，即 $\rho(x)=-\mathrm{i}\nu\frac{\delta}{\delta\phi(x)}$。利用这一关系，可以证明满足式 (7.4.6) 的算符由下式给出（Wen, 1992）：
\begin{equation*}
\Psi\propto\mathrm{e}^{\mathrm{i}\frac{1}{\nu}\phi}
\tag{7.4.7}
\end{equation*}

式 (7.4.6) 只表明算符 $\Psi$ 携带电荷 $e$。为了把 $\Psi$ 认作电子算符，我们还需要证明 $\Psi$ 是费米性算符。利用 K--M 代数 (7.4.5)，我们发现（见问题 7.4.8）
\begin{equation*}
\Psi(x)\Psi(x')=(-)^{1/\nu}\Psi(x')\Psi(x)
\tag{7.4.8}
\end{equation*}

我们看到，式 (7.4.7) 中的电子算符 $\Psi$ 只有当 $1/\nu=m$ 为奇整数时才是费米性的，此时该 FQH 态是一个 Laughlin 态。

在上述讨论中，我们做了一个并不普遍成立的假设。我们假设了不可压缩 FQH 液体只包含\emph{一种}不可压缩流体组分，从而只导致一支边缘激发。上述结果意味着，当 $\nu\neq1/m$ 时，只有一支的边缘理论不包含电子算符，因而是不自洽的。我们的单支边缘理论只适用于 $\nu=1/m$ 的 Laughlin 态。

现在让我们计算 $\nu=1/m$ Laughlin 态边缘上的电子传播子。由于 $\phi$ 是一个自由声子场，其传播子为
\begin{equation*}
\langle\phi(x,t)\phi(0)\rangle=-\nu\ln(x-vt)+\text{常数}
\end{equation*}
电子传播子可以容易地算出为
\begin{equation*}
G(x,t)=\langle T(\Psi^\dagger(x,t)\Psi(0))\rangle=\exp[\frac{1}{\nu^2}\langle\phi(x,t)\phi(0)\rangle]\propto\frac{1}{(x-vt)^m}
\tag{7.4.9}
\end{equation*}

我们看到的第一件事是：FQH 态边缘上的电子传播子获得了一个不等于 1 的非平庸指数 $m=1/\nu$。这意味着 FQH 态边缘上的电子是强关联的，不能用费米液体理论来描述。我们将把这类电子态称为手征 Luttinger 液体（chiral Luttinger liquid）。

我们想强调，指数 $m$ 是量子化的。指数的量子化与下述事实直接相关：该指数与电子的统计性质联系在一起（见式 (7.4.8)）。因此，该指数是一个拓扑数，与电子相互作用、边缘势等无关。尽管该指数是边缘态的性质，改变它的唯一途径却是通过体态中的相变。因此，这个指数可以看作一个表征\emph{体内} FQH 态拓扑序的量子数。

在动量空间中，电子传播子具有形式
\begin{equation*}
G(k,\omega)\propto\frac{(vk+\omega)^{m-1}}{\omega-vk+\mathrm{i}0^+\operatorname{sgn}(\omega)}
\end{equation*}
这个反常指数 $m$ 可以在隧穿实验中测量。由 $\operatorname{Im}G$ 可得，电子的隧穿态密度为
\begin{equation*}
N(\omega)\propto|\omega|^{m-1}
\end{equation*}
这表明，对金属--绝缘体--FQH 结，微分电导具有 $\frac{\mathrm{d}I}{\mathrm{d}V}\propto V^{m-1}$ 的形式。

\subsection*{问题 7.4.4}

证明式 (7.4.8)。（提示：先证明 $[\phi(x),\phi(y)]=\frac{\pi}{q}\operatorname{sgn}(x-y)$；译注：原文之 $q$ 疑为 $\nu$ 之误。）

\subsection*{问题 7.4.5}

\textbf{玻色化与费米化（bosonization and fermionization）}当 $m=1$ 时，本节的讨论意味着 $\nu=1$ 的 IQH 边缘态可以用自由声子理论
\begin{equation*}
H_b=2\pi v\sum_{k>0}\rho_k^\dagger\rho_k
\end{equation*}
来描述。在 7.4.1 节中，我们发现 $\nu=1$ 的 IQH 边缘态可以用自由手征费米子模型
\begin{equation*}
H_f=\sum_k vk:c_k^\dagger c_k:=\sum_{k>0}vkc_k^\dagger c_k-\sum_{k<0}vkc_kc_k^\dagger
\end{equation*}
来描述。正规排序定义为：当 $k>0$ 时 $:c_k^\dagger c_k:=c_k^\dagger c_k$；当 $k<0$ 时 $:c_k^\dagger c_k:=-c_kc_k^\dagger$。在本问题中，我们将研究这两种描述之间的直接关系。从费米子描述切换到玻色子描述称为玻色化，从玻色子描述切换到费米子描述称为费米化。
\begin{enumerate}
\item 求玻色模型前五个能级的能量及其简并度。将你的结果与问题 7.4.2 的结果作比较。
\item 令 $\rho_f(x)=:c^\dagger(x)c(x):$。在 $k$ 空间中，$\rho_{f,k}=\sum_q{}^{\prime}:c_q^\dagger c_{q+k}:$，其中求和 $\sum_q{}^{\prime}$ 限制在使 $c_{q+k}$ 与 $c_q$ 的动量都处于 $[-\Lambda,+\Lambda]$ 范围之内的那些项。正规排序后的 $\rho_f(x)$ 在基态中的平均值为零，这与玻色理论中的 $\rho$ 一致。证明 $\rho_{f,k}$ 满足与 $\rho_k$ 相同的 K--M 代数（见式 (7.4.5)，取 $\nu=1$）。（提示：对大的正 $k$ 有 $c_k^\dagger c_k=0$，对大的负 $k$ 也有 $c_k^\dagger c_k=0$。）
\item 证明 $H_f=2\pi v\sum_{k>0}\rho_{f,k}^\dagger\rho_{f,k}$。
\end{enumerate}

\subsection{边缘激发的微观理论}

\begin{keypoints}
\item K--M 代数与 Laughlin 波函数之间的关系。
\item 由等离子体类比得到的电子等时关联。
\end{keypoints}

在这一节中，我们将给出 Laughlin 态边缘激发的一个微观理论。具体起见，让我们考虑第一朗道能级中的电子气体。我们假设电子由理想哈密顿量 (7.2.5) 描述。$\nu=1/3$ 的 Laughlin 波函数
\begin{equation*}
\Psi_3(z_i)=Z^{-1/2}\prod_{i<j}(z_i-z_j)^3\prod_k\mathrm{e}^{-|z_k|^2/4l_B^2}
\tag{7.4.10}
\end{equation*}
具有零能量\footnote{这里，为方便起见，我们把能量零点移到了第零朗道能级。}，并且是理想哈密顿量的一个精确基态。（译注：原书指数中印作 $4l_B$，应为 $4l_B^2$。）在这一节中，我们将假设磁长度 $l_B=1$。式 (7.4.10) 中的 $Z$ 是归一化因子。然而，Laughlin 态 (7.4.10) 并不是唯一的零能量态。容易检验，下列类型的态都具有零能量：
\begin{equation*}
\Psi(z_i)=P(z_i)\Psi_3(z_i)
\tag{7.4.11}
\end{equation*}
其中 $P(z_i)$ 是 $z_i$ 的对称多项式。事实上，反之亦然：所有零能态都具有 (7.4.11) 的形式。这是因为，要使一个费米子态具有零能量，当任意两个电子 $i$ 与 $j$ 相互靠拢时，$\Psi$ 必须至少以 $(z_i-z_j)^3$ 的快慢消失（$(z_i-z_j)^2$ 的可能性被费米统计所排除）。由于 Laughlin 波函数只在 $z_i=z_j$ 时为零，故 $P=\Psi/\Psi_3$ 是一个有限函数。由于 $\Psi$ 与 $\Psi_3$ 都是第一朗道能级中的反对称函数，$P$ 是一个对称的全纯函数，因而只能是一个对称多项式。

在式 (7.4.11) 的所有态中，Laughlin 态描述半径最小的一个圆形液滴。所有其他的态都是 Laughlin 态液滴的形变和/或膨胀。因此，由 $P$ 生成的态对应于 Laughlin 态的边缘激发。

首先，让我们考虑零能空间（即对称多项式的空间）。我们知道，对称多项式的空间由多项式 $s_n=\sum_iz_i^n$（通过乘法与加法）生成。令 $M_0=3\frac{N(N-1)}{2}$ 为 Laughlin 态 (7.4.10) 的总角动量，则态 $\Psi$ 将具有角动量 $M=\Delta M+M_0$，其中 $\Delta M$ 是对称多项式 $P$ 的次数。由于我们只有一个零次对称多项式和一个一次对称多项式，即 $s_0=1$ 与 $s_1=\sum_iz_i$，故 $\Delta M=0,1$ 的零能态是非简并的。然而，当 $\Delta M=2$ 时，我们有两个由 $P=s_2$ 与 $P=s_1^2$ 生成的零能态。对一般的 $\Delta M$，零能态的简并度由下式给出：
\begin{equation*}
\begin{array}{cccccccc}
\Delta M: & 0 & 1 & 2 & 3 & 4 & 5 & 6\\
\text{简并度}: & 1 & 1 & 2 & 3 & 5 & 7 & 11
\end{array}
\tag{7.4.12}
\end{equation*}

这里我们想指出，式 (7.4.12) 中的简并度正是我们从宏观理论所预期的。我们知道，对于圆形液滴，角动量 $\Delta M$ 可以视为沿边缘的动量 $k=2\pi\Delta M/L$，其中 $L$ 是 FQH 液滴的周长。按照宏观理论，（中性的）边缘激发由密度算符 $\rho_k$ 生成。容易检验，由密度算符生成的边缘态对每一个 $\Delta M$ 都具有与式 (7.4.12) 相同的简并度。例如，$\Delta M=2$ 的两个态由 $\rho_{\kappa_0}^2$ 与 $\rho_{2\kappa_0}$ 生成，其中 $\kappa_0=2\pi/L$。因此，由 K--M 代数 (7.4.5) 生成的空间与对称多项式的空间是同一个空间。

现在让我们问下面这个物理问题：对称多项式是否生成所有的低能态？如果是这样，那么由上述讨论我们看到，FQH 液滴（译注：原文印作“HQ 液滴”）的所有低能激发都由 K--M 代数生成，我们就可以说式 (7.4.5) 是低能激发的一个完备理论。遗憾的是，到目前为止我们还没有上述论断的解析证明。这是因为，虽然与对称多项式所生成的态正交的态具有非零能量，但这些能量在热力学极限下是否保持有限并不清楚。能隙有可能在热力学极限下趋于零。为了解决这个问题，目前我们只能依靠数值计算。在图 7.17 中，我们给出了六个电子的系统在前二十二个轨道上、对本节开头所引入哈密顿量的能谱。零能态在 $M=45,\ldots,51$（即 $\Delta M=0,\ldots,6$）处的简并度分别求得为 $1,1,2,3,5,7,11$，与式 (7.4.12) 相符。更重要的是，我们清楚地看到，一个有限的能隙把所有其余的态与零能态分隔开。因此，数值结果意味着，Laughlin 态的所有低能边缘激发都由对称多项式或 K--M 代数 (7.4.5) 生成。

%FIG{7.17}: {六个电子的系统在前二十二个轨道中对于哈密顿量 $H_V$ 的能谱。$M=45,\ldots,51$ 处零能态的简并度分别求得为 $1,1,2,3,5,7,11$。}

在下面，我们将从 Laughlin 波函数 $\Psi_m$ 出发计算准粒子/电子的等时关联。首先，我们如下计算 $|\xi\rangle=\prod_i(1-\frac{z_i}{\xi})^n\Psi_m(\{z_i\})$ 的模：
\begin{align*}
\langle\xi|\xi\rangle&=|\xi|^{-2Nn}\frac{Z_1}{Z}\tag{7.4.13}\\
Z&=\int\prod\mathrm{d}^2z_i\ \exp\Bigl(-\beta\sum_{ij}(-m^2)\ln|z_i-z_j|-\beta\sum_k\frac{m}{4}|z_k|^2\Bigr)\\
Z_1&=\int\prod\mathrm{d}^2z_i\ \exp\Bigl[\beta\sum_{ij}m^2\ln|z_i-z_j|-\beta\sum_k\Bigl(\frac{m}{4}|z_k|^2-nm\ln|\xi-z_k|\Bigr)\Bigr]
\end{align*}
其中 $\beta=2/m$。注意，$Z$ 是由 $N$ 个“电荷”为 $m$ 的粒子组成的等离子体的配分函数，而 $Z_1$ 是该等离子体与 $\xi$ 处一个“电荷”为 $n$ 的粒子发生相互作用的配分函数。我们可以写成
\begin{equation*}
Z=\mathrm{e}^{-\beta E},\qquad Z_1=\mathrm{e}^{-\beta E_1}
\end{equation*}
其中 $E$ 与 $E_1$ 是等离子体的总能量。能量的改变 $E_1-E$ 由加在 $\xi$ 处的“电荷” $n$ 粒子与 $z_i$ 处 $N$ 个“电荷” $m$ 粒子之间的相互作用给出。当 $|\xi|$ 大时，$N$ 个“电荷” $m$ 粒子形成一个圆形液滴，我们可以把它们当作 $z=0$ 处的一个点“电荷” $mN$ 来处理。我们得到 $E_1-E=-nmN\ln|\xi|$。对于较小的 $|\xi|$，“电荷” $n$ 会改变液滴的形状，此时
\begin{equation*}
E_1-E=-nmN\ln|\xi|-\frac{n^2}{2}\Bigl[-\ln\Bigl(|\xi|-\frac{R^2}{|\xi|}\Bigr)+\ln(|\xi|)\Bigr]
\tag{7.4.14}
\end{equation*}
其中 $R=\sqrt{mN}\,l_B$ 是液滴的半径。式 (7.4.14) 中的第一项是“电荷” $n$ 与\emph{未形变}液滴之间的相互作用。第二项是液滴形变引起的修正。由于等离子体的行为像金属，这一修正可以用“电荷” $n$ 与它在 $|z|=R^2/|\xi|$ 及 $z=0$ 处的镜像之间的相互作用来表示。

在上述讨论中，我们忽略了电荷的分立性，把等离子体当作连续介质来处理。我们预期，只要 $\xi$ 不太靠近液滴，即 $|\xi|-R\gg l_B$，这一近似就能给出正确的比值 $Z_1/Z$。

由式 (7.4.14)，我们求得 $|\psi^n(\xi)\rangle$ 的模为
\begin{equation*}
\langle\xi|\xi\rangle=\Bigl(\frac{\xi\xi^*}{\xi\xi^*-R^2}\Bigr)^{n^2/m}
\tag{7.4.15}
\end{equation*}
由于内积 $\langle\tilde{\xi}|\xi\rangle$ 是 $\xi$ 的全纯函数、$\tilde{\xi}$ 的反全纯函数，式 (7.4.15) 意味着
\begin{equation*}
\langle\tilde{\xi}|\xi\rangle=\Bigl(\frac{\xi\tilde{\xi}^*}{\xi\tilde{\xi}^*-R^2}\Bigr)^{n^2/m}
\end{equation*}
注意，当 $n=m$ 时，
\begin{equation*}
\xi^{Nm}\tilde{\xi}^{*Nm}\langle\tilde{\xi}|\xi\rangle=\xi^{Nm}\tilde{\xi}^{*Nm}\Bigl(\frac{\xi\tilde{\xi}^*}{\xi\tilde{\xi}^*-R^2}\Bigr)^m
\end{equation*}
正比于 $N_e=N+1$ 电子系统边缘上等时的电子传播子 $G_e$。取 $\xi=R\mathrm{e}^{\mathrm{i}\frac{x}{R}}$ 与 $\tilde{\xi}=R$，我们得到
\begin{equation*}
G_e(x)=L^{-m}a^{m-1}\mathrm{e}^{\mathrm{i}m\left(N_e-\frac{1}{2}\right)\frac{2\pi x}{L}}\sin^{-m}(\pi x/L)
\tag{7.4.16}
\end{equation*}
当 $x$ 远小于 $L\equiv2\pi R$（译注：原文印作“$21R$”，应为液滴周长 $2\pi R$）时，上式退化为式 (7.4.9)。这里 $a$ 是 $l_B$ 量级的一个长度标度。式 (7.4.16) 可以展开如下：
\begin{align*}
G_e(x)&=L^{-m}a^{m-1}\mathrm{e}^{\mathrm{i}m(N_e-1)\frac{2\pi x}{L}}\Bigl[\sum_{n=0}^{\infty}\mathrm{e}^{-\mathrm{i}\frac{2\pi x}{L}n}\Bigr]^m\notag\\
&=L^{-m}a^{m-1}\mathrm{e}^{\mathrm{i}m(N_e-1)\frac{2\pi x}{L}}\sum_{n=0}^{\infty}C_n^{m+n-1}\,\mathrm{e}^{-\mathrm{i}\frac{2\pi x}{L}n}\\
C_n^{m+n-1}&=\frac{(n+m-1)!}{(m-1)!\,n!}\tag{7.4.17}
\end{align*}

从这个展开，我们得到角动量 $M$ 态上的电子占据数 $n_M$ 如下：
\begin{equation*}
\begin{array}{ll}
n_M=0, & \delta M=m(N_e-1)-M<0\\[4pt]
n_M=\dfrac{a^{m-1}}{L^{m-1}}C_{\delta M}^{\delta M+m-1}, & \delta M\geqslant0
\end{array}
\tag{7.4.18}
\end{equation*}
我们看到，费米边缘的准确位置处在最后一个部分占据的单粒子轨道上，即处在角动量 $m(N_e-1)$ 处（或 $k_F=\sqrt{m(N_e-1)}/l_B$ 处）。注意，当 $m(N_e-1)-M\gg m$ 时，我们有 $n_M\propto(m(N_e-1)-M)^{m-1}$。用沿边缘的动量来表示，我们有 $n_k\propto(k_F-k)^{m-1}$。我们发现，与费米液体不同，占据数 $n_k$ 在费米动量处没有跳跃。

我们想指出，式 (7.4.16) 只在 $x$ 远大于磁长度 $l_B$ 时才正确。因此，式 (7.4.18) 只在 $m(N_e-1)-M\ll\sqrt{N_e}$ 时才有效。

如果边缘激发的色散是线性的，那么在低能下 $G_e$ 只能依赖于 $x-vt$。我们立刻看到，含时的电子格林函数为
\begin{equation*}
G_e(x,t)=L^{-m}a^{m-1}\mathrm{e}^{\mathrm{i}m\left(N-\frac{1}{2}\right)\frac{2\pi(x-vt)}{L}}\sin^{-m}[\pi(x-vt)/L]
\end{equation*}
当 $m=1$ 时，上式成为一维自由费米子的电子格林函数（对费米点附近的 $k$）。

\subsection{流体动力学方法——2/5 与 2/3 态}

\begin{keypoints}
\item 分级 FQH 液体的边缘态。
\item 电子算符与准粒子算符的结构。
\item 当且仅当边缘激发能够沿相反方向传播时，边缘相互作用才能修改电子/准粒子传播子的指数。
\item 当 $\nu_1=-\nu_2=1$ 时，本节的讨论描述一维 Tomonaga--Luttinger 液体。
\end{keypoints}

在这一节中，我们将用流体动力学方法研究第二级分级态的边缘结构。我们将以 2/5 态与 2/3 态为例。特别地，我们将研究分级态边缘上电子算符与准粒子算符的结构。我们还将看到，2/3 态包含两支沿相反方向传播的边缘模式，这是相当反直觉的。

首先，让我们考虑 $\nu=\frac{2}{5}$ 的 FQH 态。按照分级理论，$\nu=\frac{2}{5}$ 的 FQH 态是在 $\nu=\frac{1}{3}$ 的 FQH 态之上凝聚准粒子而生成的。因此，2/5 态包含两种不可压缩流体组分。为了明确起见，让我们考虑一个特殊的边缘势，使得 FQH 态由两个液滴组成（见图 7.14）：一个是电子凝聚体，填充因子为 $\nu_1=\frac{1}{3}$，半径为 $r_1$；另一个是准粒子凝聚体（位于 1/3 态之上），填充因子为 $\nu_2=\frac{1}{15}$（注意 $\frac{1}{3}+\frac{1}{15}=\frac{2}{5}$），半径为 $r_2<r_1$。

当 $r_1-r_2\gg l_B$ 时，两条边缘相互独立。推广 7.4.2 节的流体动力学方法，可以证明存在两支边缘激发，其低能动力学由下式描述：
\begin{equation*}
\begin{gathered}
[\rho_{Ik},\rho_{Jk'}]=\frac{\nu_I}{2\pi}k\delta_{IJ}\delta_{k+k'}\\
H=2\pi\sum_{I,k>0}\frac{v_I}{\nu_I}\rho_{I,-k}\rho_{I,k}
\end{gathered}
\tag{7.4.19}
\end{equation*}
其中 $I=1,2$ 标记两支，$v_I$ 是边缘激发的速度。在式 (7.4.19) 中，$\rho_I$ 是一维电子密度。为了使哈密顿量下有界，我们要求 $\nu_Iv_I>0$。我们发现，$\nu=\frac{2}{5}$ FQH 态的稳定性要求两个 $v_I$ 都为正。

推广 7.4.2 节的讨论，两条边缘上的电子算符为
\begin{equation*}
\Psi_I=\mathrm{e}^{\mathrm{i}\frac{1}{\nu_I}\phi_I(x)},\qquad I=1,2
\tag{7.4.20}
\end{equation*}
其中 $\partial_x\phi_I=\frac{1}{2\pi}\rho_I$。电子传播子具有形式
\begin{equation*}
\langle T(\Psi_I(x,t)\Psi_I^\dagger(0))\rangle=\frac{\mathrm{e}^{\mathrm{i}k_Ix}}{(x-v_I t)^{-1/|\nu_I|}},\qquad I=1,2
\end{equation*}
这里 $k_I=r_I/2l_B^2$。

按照分级图像，$\nu=\frac{2}{3}$ 的 FQH 态也是由两个凝聚体形成的：一个填充因子为 1 的电子凝聚体和一个填充因子为 $-\frac{1}{3}$ 的空穴凝聚体。因此，取 $(\nu_1,\nu_2)=(1,-\frac{1}{3})$，上述讨论同样可以应用于 $\nu=\frac{2}{3}$ 的 FQH 态。同样存在两支边缘激发，但只要哈密顿量是正定的，这两支现在就具有\emph{相反}的速度。

当我们把两条边缘靠拢到 $r_1-r_2\sim l_B$ 时，两支边缘激发之间的相互作用不再可以忽略。此时，哈密顿量具有形式
\begin{equation*}
H=2\pi\sum_{k>0}V_{IJ}\rho_{I,-k}\rho_{J,k}
\tag{7.4.21}
\end{equation*}

哈密顿量 (7.4.21) 可以对角化。当 $\nu_1\nu_2>0$ 时，我们可以选择
\begin{align*}
\tilde{\rho}_{1k}&=\cos(\theta)\frac{1}{\sqrt{|\nu_1|}}\rho_{1k}+\sin(\theta)\frac{1}{\sqrt{|\nu_2|}}\rho_{2k}\\
\tilde{\rho}_{2k}&=\cos(\theta)\frac{1}{\sqrt{|\nu_2|}}\rho_{2k}-\sin(\theta)\frac{1}{\sqrt{|\nu_1|}}\rho_{1k}\\
\tan(2\theta)&=2\frac{\sqrt{|\nu_1\nu_2|}V_{12}}{|\nu_1|V_{11}-|\nu_2|V_{22}}
\tag{7.4.22}
\end{align*}
可以检验，$\tilde{\rho}$ 满足
\begin{equation*}
\begin{gathered}
[\tilde{\rho}_{Ik},\tilde{\rho}_{Jk'}]=\frac{\operatorname{sgn}(\nu_I)}{2\pi}k\delta_{IJ}\delta_{k+k'}\\
H=2\pi\sum_{I,k>0}\operatorname{sgn}(\nu_I)\tilde{v}_I\tilde{\rho}_{I,-k}\tilde{\rho}_{I,k}
\end{gathered}
\tag{7.4.23}
\end{equation*}
其中边缘激发的新速度 $\tilde{v}_I$ 由下式给出：
\begin{align*}
\operatorname{sgn}(\nu_1)\tilde{v}_1&=\frac{\cos^2(\theta)}{\cos(2\theta)}|\nu_1|V_{11}-\frac{\sin^2(\theta)}{\cos(2\theta)}|\nu_2|V_{22}\\
\operatorname{sgn}(\nu_2)\tilde{v}_2&=\frac{\cos^2(\theta)}{\cos(2\theta)}|\nu_2|V_{22}-\frac{\sin^2(\theta)}{\cos(2\theta)}|\nu_1|V_{11}
\tag{7.4.24}
\end{align*}
我们看到仍然有两支边缘激发。然而在这种情况下，具有确定速度的边缘激发是内边缘上的激发与外边缘上的激发的混合。还可以证明，只要哈密顿量 (7.4.21) 下有界，两支的速度 $\tilde{v}_I$ 就总是正的。

通过反解式 (7.4.22) 把式 (7.4.20) 中的电子算符 $\Psi_I$ 用 $\tilde{\rho}_I$ 重写之后，我们可以利用式 (7.4.24) 计算它们的传播子如下：
\begin{equation*}
\langle T(\Psi_I(x,t)\Psi_I^\dagger(0))\rangle=\mathrm{e}^{\mathrm{i}k_Ix}\frac{1}{(x-\tilde{v}_1t)^{\alpha_I}}\frac{1}{(x-\tilde{v}_2t)^{\beta_I}}
\end{equation*}
其中
\begin{equation*}
(\alpha_1,\alpha_2)=\Bigl(\frac{\cos^2\theta}{|\nu_1|},\frac{\sin^2\theta}{|\nu_2|}\Bigr),\qquad
(\beta_1,\beta_2)=\Bigl(\frac{\sin^2\theta}{|\nu_1|},\frac{\cos^2\theta}{|\nu_2|}\Bigr)
\end{equation*}

然而，当两条边缘靠近到磁长度以内时，$\Psi_I$ 就不再是边缘上最普遍的电子算符了。普遍的电子算符可以包含两边缘之间的电荷转移。对 $\nu=2/5$ 的 FQH 态，内边缘与外边缘之间隔着 $\nu=\frac{1}{3}$ 的 Laughlin 态。因此，基本的电荷转移算符由下式给出：
\begin{equation*}
\eta(x)=\mathrm{e}^{\mathrm{i}(\phi_1-\frac{\nu_1}{\nu_2}\phi_2)}=(\Psi_1\Psi_2^\dagger)^{\nu_1}
\end{equation*}
它把一份 $\nu_1e=e/3$ 的电荷从外边缘转移到内边缘。普遍的电子算符于是具有形式
\begin{equation*}
\begin{gathered}
\Psi(x)=\sum_{n=-\infty}^{+\infty}c_n\psi_n(x)\\
\psi_n(x)=\Psi_1(x)\eta^n(x)
\end{gathered}
\tag{7.4.25}
\end{equation*}

为了理解这个结果，我们注意到，无论整数 $n$ 取何值，每个算符 $\psi_n$ 总是产生一个单位的局域电荷，并且是费米性算符。因此，每个 $\psi_n$ 都是边缘电子算符的一个候选者。对一个普遍的相互作用系统，边缘上的电子算符预期是不同 $\psi_n$ 的叠加，如式 (7.4.25) 所表示。注意 $\Psi_2=\psi_{-1/\nu_1}$。$\psi_n$ 的传播子可以用与上面概述类似的方法计算，结果为
\begin{equation*}
\langle T(\psi_n(x,t)\psi_m^\dagger(0))\rangle\propto\delta_{n,m}\,\mathrm{e}^{\mathrm{i}[k_1+n\nu_1(k_2-k_1)]x}\prod_I(x-\tilde{v}_It)^{-\gamma_{In}}
\end{equation*}
其中 $\gamma_{In}$ 为
\begin{align*}
\gamma_{1n}&=\Bigl[\Bigl(n+\frac{1}{|\nu_1|}\Bigr)\sqrt{|\nu_1|}\cos\theta-\frac{n\nu_1}{\nu_2}\sqrt{|\nu_2|}\sin\theta\Bigr]^2\\
\gamma_{2n}&=\Bigl[\Bigl(n+\frac{1}{|\nu_1|}\Bigr)\sqrt{|\nu_1|}\sin\theta+\frac{n\nu_1}{\nu_2}\sqrt{|\nu_2|}\cos\theta\Bigr]^2
\tag{7.4.26}
\end{align*}

从式 (7.4.25) 与上式（译注：原文作“式 (7.4.4)”，应指上面未编号的 $\psi_n$ 传播子公式）我们看到，电子传播子在分立动量 $k=k_1+n\nu_1(k_2-k_1)$ 处具有奇异性。它们类似于相互作用一维电子系统中电子传播子的 $k_F,3k_F,\ldots$ 奇异性。

对 $\nu=\frac{2}{3}$ 的 FQH 态，我们有 $\nu_1\nu_2<0$。在这种情况下，我们需要选择
\begin{align*}
\tilde{\rho}_{1k}&=\cosh(\theta)\frac{1}{\sqrt{|\nu_1|}}\rho_{1k}+\sinh(\theta)\frac{1}{\sqrt{|\nu_2|}}\rho_{2k}\\
\tilde{\rho}_{2k}&=\cosh(\theta)\frac{1}{\sqrt{|\nu_2|}}\rho_{2k}+\sinh(\theta)\frac{1}{\sqrt{|\nu_1|}}\rho_{1k}\\
\tanh(2\theta)&=2\frac{\sqrt{|\nu_1\nu_2|}V_{12}}{|\nu_1|V_{11}+|\nu_2|V_{22}}
\tag{7.4.27}
\end{align*}
来对角化哈密顿量。可以检验，$\tilde{\rho}_I$ 同样满足 K--M 代数 (7.4.23)，只是现在
\begin{align*}
\operatorname{sgn}(\nu_1)\tilde{v}_1&=\frac{\cosh^2(\theta)}{\cosh(2\theta)}|\nu_1|V_{11}-\frac{\sinh^2(\theta)}{\cosh(2\theta)}|\nu_2|V_{22}\\
\operatorname{sgn}(\nu_2)\tilde{v}_2&=\frac{\cosh^2(\theta)}{\cosh(2\theta)}|\nu_2|V_{22}-\frac{\sinh^2(\theta)}{\cosh(2\theta)}|\nu_1|V_{11}
\tag{7.4.28}
\end{align*}
同样，只要哈密顿量 $H$ 是正定的，边缘激发的速度 $\tilde{v}_I$ 就总是具有相反的符号。电子算符仍具有 (7.4.25) 的形式，其中 $\eta=(\Psi_1\Psi_2^\dagger)^{\nu_1}$。$\psi_n$ 的传播子仍由式 (7.4.4)（译注：同前，应指上面未编号的 $\psi_n$ 传播子公式）给出，其中
\begin{align*}
\gamma_{1n}&=\Bigl[\Bigl(n+\frac{1}{|\nu_1|}\Bigr)\sqrt{|\nu_1|}\cosh\theta+\frac{n\nu_1}{\nu_2}\sqrt{|\nu_2|}\sinh\theta\Bigr]^2\\
\gamma_{2n}&=\Bigl[\Bigl(n+\frac{1}{|\nu_1|}\Bigr)\sqrt{|\nu_1|}\sinh\theta+\frac{n\nu_1}{\nu_2}\sqrt{|\nu_2|}\cosh\theta\Bigr]^2
\tag{7.4.29}
\end{align*}

在等空间点处，电子算符 $\psi_n$ 的长时间关联由总指数 $g_e^{(n)}$ 控制如下：
\begin{equation*}
\langle\psi_n^\dagger(0,t)\psi_n(0,0)\rangle\sim t^{-g_e^{(n)}},\qquad g_e^{(n)}=\sum_I\gamma_{In}
\end{equation*}

等空间的电子关联决定了边缘隧穿实验中的 $I$--$V$ 曲线（见 3.7.4 节与 4.2.2 节）。指数的最小值 $g_e\equiv\operatorname{Min}(g_e^{(n)})$ 控制着电子在两边缘之间隧穿的标度性质。例如，隧穿电导在有限温度下的标度行为为
\begin{equation*}
\sigma\propto T^{2g_e-2}
\end{equation*}
对 $\nu=2/5$ 态，这里 $g_e=3$。对 $\nu=2/3$ 态，$g_e$ 依赖于两边缘之间的相互作用。

\subsection*{问题 7.4.6}

\textbf{一维相互作用费米子系统的玻色化}考虑一个自由的一维费米子模型
\begin{equation*}
H_0=\sum_k\epsilon_k:c_k^\dagger c_k:
\end{equation*}
正规排序定义为：当 $\epsilon_k>0$ 时 $:c_k^\dagger c_k:=c_k^\dagger c_k$；当 $\epsilon_k<0$ 时 $:c_k^\dagger c_k:=-c_kc_k^\dagger$。这里当 $k=\pm k_F$ 时 $\epsilon_k=0$。
\begin{enumerate}
\item 令 $\rho_{1,k}=\sum_{q\sim k_F}^{\prime}:c_q^\dagger c_{q+k}:$，其中求和 $\sum_{q\sim k_F}^{\prime}$ 限制在使 $c_{q+k}$ 与 $c_q$ 的动量都处于范围 $[-\Lambda+k_F,+\Lambda+k_F]$ 之内的那些项。令 $\rho_{2,k}=\sum_{q\sim-k_F}^{\prime}:c_q^\dagger c_{q+k}:$，其中求和 $\sum_{q\sim-k_F}^{\prime}$ 使 $c_{q+k}$ 与 $c_q$ 处于范围 $[-\Lambda-k_F,+\Lambda-k_F]$ 之内。这里 $k<\Lambda<k_F$。我们注意到，$\rho_{1,k}$ 在 $k_F$ 附近产生右行激发，而 $\rho_{2,k}$ 在 $-k_F$ 附近产生左行激发。证明 $\rho_{I,k}$ 满足 K--M 代数 (7.4.19)，其中 $\nu_1=1$，$\nu_2=-1$。证明 $H_0=2\pi v_F\sum_{k>0}(\rho_{1,k}\rho_{1,-k}-\rho_{2,-k}\rho_{2,k})$。于是一维自由费米子系统可以被玻色化。
\item 我们可以把电子算符写成 $c(x)=\psi_1(x)+\psi_2(x)$，其中 $\psi_1(x)=\sum_{-\Lambda+k_F}^{\Lambda+k_F}\mathrm{e}^{\mathrm{i}kx}c_k$ 是 $k_F$ 附近的电子算符，而 $\psi_2(x)=\sum_{-\Lambda-k_F}^{\Lambda-k_F}\mathrm{e}^{\mathrm{i}kx}c_k$ 是 $-k_F$ 附近的电子算符。试用 $\phi_I$ 表示 $c(x)$，这里 $(2\pi)^{-1}\partial_x\phi_I(x)=\rho_I(x)$。
\item 为了验证一维相互作用费米子系统也可以被玻色化，证明相互作用项 $H_V=\sum_{q,k_1,k_2}V_q\allowbreak(c_{k_1}^\dagger c_{k_1+q})\allowbreak(c_{k_2}^\dagger c_{k_2+q})$ 在限制到两个费米点附近时变为 $2\pi\sum\allowbreak V_{IJ}\rho_{I,-k}\rho_{J,k}+$ 常数。求 $V_{IJ}$。
\item 计算相互作用模型 $H_0+H_V$ 的电子格林函数 $\langle c^\dagger(x,t)c(0)\rangle$。
\end{enumerate}

\subsection{体有效理论与边缘态}

\begin{keypoints}
\item 边缘态的结构可以由以 $K$ 矩阵表征的体拓扑序直接确定。
\end{keypoints}

在这一节中，我们将从体 FQH 态的 Chern--Simons 有效理论直接导出边缘激发的宏观理论。在这种方法中，我们不依赖于 FQH 态的任何具体构造。体拓扑序与边缘态之间的关系变得十分清楚。我必须提醒读者，本节所给出的计算是非常形式的。结果的正确性不是由形式计算来保证的，而是由与其他独立计算（例如 7.4.3 节中的那个）的比较来保证的。

为了理解有效理论与边缘态之间的关系，让我们先考虑填充因子 $\nu=1/m$ 的最简单 FQH 态，并尝试由体有效理论重新导出 7.4.2 节的结果。这样的 FQH 态由具有作用量
\begin{equation*}
S=-\frac{m}{4\pi}\int a_\mu\partial_\nu a_\lambda\varepsilon^{\mu\nu\lambda}\mathrm{d}^3x
\tag{7.4.30}
\end{equation*}
的 U(1) Chern--Simons 理论描述。

假设我们的样品具有一个边界。为简单起见，我们假设边界是 $x$ 轴，样品覆盖下半平面。

对于带边界的 FQH 液体，有效作用量 (7.4.30) 存在一个问题：由于边界的存在，它在规范变换 $a_\mu\to a_\mu+\partial_\mu f$ 下不是不变的，即 $\Delta S=-\frac{m}{4\pi}\int_{y=0}\mathrm{d}x\mathrm{d}t\,f(\partial_ta_1-\partial_1a_0)$。为了解决这个问题，我们将把规范变换限制为在边界上为零：$f(x,y=0,t)=0$。由于这一限制，$a_\mu$ 在边界上的一些自由度变成了动力学自由度。

我们知道，有效理论 (7.4.30) 只是针对无边界的体 FQH 态导出的。这里，我们将把带有受限规范变换的式 (7.4.30) 取作为带边界的 FQH 态的有效理论的定义。这样的定义肯定是自洽的。下面我们将证明，这样的定义重现了 7.4.2 节所得到的结果。

研究规范理论动力学的一种方法是选择规范条件 $a_0=0$，并把 $a_0$ 的运动方程当作一个约束。对 Chern--Simons 理论，这样的约束成为 $f_{ij}=0$。在这个约束下，我们可以把 $a_i$ 写成 $a_i=\partial_i\phi$。把它代入式 (7.4.30)，就得到边缘上一个有效的一维理论，其作用量为（Elitzur et al., 1989）
\begin{equation*}
S_{\text{edge}}=-\frac{m}{4\pi}\int\partial_t\phi\,\partial_x\phi\,\mathrm{d}x\mathrm{d}t
\tag{7.4.31}
\end{equation*}

然而，这个方法有一个问题。容易看出，与作用量 (7.4.31) 相联系的哈密顿量为零，而且式 (7.4.31) 所描述的边界激发的速度为零。因此，这个作用量不能用来描述真实 FQH 样品中任何物理的边缘激发。FQH 态中的边缘激发总是具有有限的速度。

边缘激发出现有限速度是一种边界效应。由式 (7.3.18) 定义的体有效理论不包含边缘激发速度的信息。为了从有效理论确定边缘激发的动力学，我们必须找到一种把边缘速度的信息输入进来的办法。边缘速度必须被当作体有效理论中不包含的外部参量来处理。问题在于如何把这些参量放进理论之中。

现在让我们注意到，$a_0=0$ 这个条件并不是规范固定条件的唯一选择。更普遍的规范固定条件具有形式
\begin{equation*}
a_\tau=a_0+va_x=0
\tag{7.4.32}
\end{equation*}
这里 $a_x$ 是矢量势平行于样品边界的分量，而 $v$ 是一个具有速度量纲的参量。

方便的做法是选择满足
\begin{equation*}
\tilde{x}=x-vt,\qquad \tilde{t}=t,\qquad \tilde{y}=y.
\tag{7.4.33}
\end{equation*}
的新坐标。注意，在坐标变换下规范势 $a_\mu$ 按 $\partial/\partial x^\mu$ 变换。我们发现在新坐标下，规范场的各分量
